\documentclass{article}
\usepackage{amsmath,amssymb,amsthm}
\usepackage{algorithm}
\usepackage{algpseudocode}
\usepackage{hyperref}
\newtheorem{theorem}{Theorem}
\newtheorem{lemma}{Lemma}
\newtheorem{assumption}{Assumption}
\newcommand{\E}{\mathbb{E}}
\newcommand{\norm}[1]{\left\|#1\right\|}
\newcommand{\inner}[2]{\left\langle #1,#2\right\rangle}
\begin{document}

\section*{AdaGrad for Non‑Convex Smooth Objectives (Scalar Version)}

\section*{Mathematical Formulation}
Let $f:\mathbb{R}\rightarrow\mathbb{R}$ be differentiable.  
For iteration $k\ge 0$ denote the stochastic (or deterministic) gradient by
\[
g_k:=\nabla f(w_k).
\]
Define the accumulated squared gradient
\[
G_k:=\sum_{i=0}^{k} g_i^{2},\qquad G_{-1}:=0 .
\]
Given a base learning rate $\alpha>0$, the AdaGrad update reads
\begin{equation}
\label{eq:update}
w_{k+1}=w_{k}-\frac{\alpha}{\sqrt{G_k+\epsilon}}\;g_k,
\end{equation}
where $\epsilon>0$ is a small constant to avoid division by zero.  
The algorithm keeps a running sum $A_k:=\sum_{i=0}^{k} g_i$ (used only for exposition; the update only needs $G_k$).

\section*{Pseudocode}
\begin{algorithm}[H]
\caption{Scalar AdaGrad for Non‑Convex $f$}
\begin{algorithmic}[1]
\State \textbf{Input:} initial point $w_0\in\mathbb{R}$, base step size $\alpha>0$, tolerance $\epsilon>0$, max iterations $T$.
\State $G_{-1}\gets 0$ \Comment{accumulated squared gradients}
\For{$k=0$ \textbf{to} $T-1$}
    \State Compute (stochastic) gradient $g_k \gets \nabla f(w_k)$
    \State $G_k \gets G_{k-1}+g_k^{2}$
    \State $\eta_k \gets \dfrac{\alpha}{\sqrt{G_k+\epsilon}}$
    \State $w_{k+1} \gets w_k-\eta_k\,g_k$
\EndFor
\State \textbf{Output:} $\{w_k\}_{k=0}^{T}$
\end{algorithmic}
\end{algorithm}

\section*{Dry Run Example}
Consider $f(w)=(w-3)^2$, thus $\nabla f(w)=2(w-3)$.  
Choose $w_0=0$, $\alpha=0.1$, $\epsilon=10^{-8}$, and run three iterations.

\begin{center}
\begin{tabular}{c|c|c|c|c}
$k$ & $w_k$ & $g_k=2(w_k-3)$ & $G_k$ & $w_{k+1}=w_k-\dfrac{\alpha}{\sqrt{G_k+\epsilon}}g_k$\\\hline
0 & $0$ & $-6$ & $36$ & $0-\dfrac{0.1}{\sqrt{36}}(-6)=0+0.1\cdot\frac{6}{6}=0.1$\\
1 & $0.1$ & $-5.8$ & $36+33.64=69.64$ & $0.1-\dfrac{0.1}{\sqrt{69.64}}(-5.8)=0.1+0.1\cdot\frac{5.8}{8.346}=0.1694$\\
2 & $0.1694$ & $-5.6612$ & $69.64+32.058=101.698$ &
$0.1694-\dfrac{0.1}{\sqrt{101.698}}(-5.6612)=0.1694+0.1\cdot\frac{5.6612}{10.084}=0.2354$
\end{tabular}
\end{center}
The objective values decrease:
\[
f(0)=9,\; f(0.1)=8.41,\; f(0.1694)=8.18,\; f(0.2354)=7.84.
\]

\newpage
\section*{Assumptions}
\begin{assumption}[Smoothness]\label{ass:smooth}
$f$ is $L$‑smooth, i.e.
\[
\forall u,v\in\mathbb{R}:\quad |f(u)-f(v)-\inner{\nabla f(v)}{u-v}|\le \frac{L}{2}\norm{u-v}^{2}.
\]
\end{assumption}
\begin{assumption}[Bounded Gradient Variance]\label{ass:variance}
The stochastic gradient $g_k$ satisfies
\[
\E[g_k\mid w_k]=\nabla f(w_k),\qquad 
\E\!\big[\|g_k-\nabla f(w_k)\|^{2}\mid w_k\big]\le\sigma^{2},
\]
for some $\sigma\ge0$.
\end{assumption}
\begin{assumption}[Bounded Second Moment]\label{ass:bound}
There exists $G_{\max}>0$ such that $\E[g_k^{2}]\le G_{\max}$ for all $k$.
\end{assumption}

\section*{Main Convergence Theorem}
\begin{theorem}[AdaGrad Gradient Norm Bound]\label{thm:main}
Under Assumptions~\ref{ass:smooth}–\ref{ass:bound}, for any $T\ge1$ the iterates generated by \eqref{eq:update} satisfy
\[
\frac{1}{T}\sum_{k=0}^{T-1}\E\!\big[\|\nabla f(w_k)\|^{2}\big]
\;\le\;
\frac{2(f(w_0)-f^{\star})}{\alpha\,\sqrt{T}} 
\;+\;
\frac{2L\alpha\sigma^{2}}{ \sqrt{T}},
\]
where $f^{\star}:=\inf_{w}f(w)$.
Consequently, $\displaystyle\min_{0\le k<T}\E\!\big[\|\nabla f(w_k)\|^{2}\big]=O\!\big(T^{-1/2}\big)$.
\end{theorem}

\section*{Theoretical Properties}
\begin{itemize}
\item \textbf{Stability.} The adaptive denominator $\sqrt{G_k+\epsilon}$ monotonically grows, guaranteeing that the effective step size $\eta_k$ is non‑increasing and thus prevents explosive updates.
\item \textbf{Hyper‑parameter Sensitivity.} The base learning rate $\alpha$ appears linearly in the bound; larger $\alpha$ accelerates early progress but may increase the variance term $2L\alpha\sigma^{2}/\sqrt{T}$.
\item \textbf{Comparison.} For a fixed step size $\eta$, the standard SGD bound is $O(1/\sqrt{T})$ with a constant proportional to $\eta L\sigma^{2}$. AdaGrad automatically scales $\eta_k$ down, yielding the same asymptotic rate without manual tuning of a diminishing schedule.
\end{itemize}

\section*{Discussion}
The theorem mirrors classical AdaGrad analyses for convex problems, yet the proof only uses smoothness and variance bounds, thus extending to non‑convex settings. The $O(T^{-1/2})$ rate is optimal for first‑order methods under the stated assumptions.

\newpage
\section*{Preliminaries (Definitions and Lemmas)}
\begin{lemma}[Smoothness Inequality]\label{lem:smooth}
If $f$ is $L$‑smooth, then for any $w$ and $g$,
\[
f(w-g)-f(w)\le -\inner{\nabla f(w)}{g}+\frac{L}{2}\|g\|^{2}.
\]
\end{lemma}
\begin{proof}
Apply Assumption~\ref{ass:smooth} with $u=w-g$, $v=w$ and rearrange. ∎
\end{proof}

\begin{lemma}[Bound on Accumulated Steps]\label{lem:stepbound}
For any non‑negative sequence $\{a_k\}$,
\[
\sum_{k=0}^{T-1}\frac{a_k}{\sqrt{\sum_{i=0}^{k}a_i+\epsilon}}\le 2\sqrt{\sum_{i=0}^{T-1}a_i+\epsilon}.
\]
\end{lemma}
\begin{proof}
Define $S_k:=\sum_{i=0}^{k}a_i+\epsilon$. Then $S_{k}=S_{k-1}+a_k$. Observe
\[
\frac{a_k}{\sqrt{S_k}}=2\big(\sqrt{S_k}-\sqrt{S_{k-1}}\big)-\frac{a_k^{2}}{2\,S_k^{3/2}}\le 2\big(\sqrt{S_k}-\sqrt{S_{k-1}}\big).
\]
Summing from $k=0$ to $T-1$ telescopes to $2(\sqrt{S_{T-1}}-\sqrt{S_{-1}})\le2\sqrt{S_{T-1}}$. ∎
\end{proof}

\section*{Main Proof (Step‑by‑Step Derivation)}
We prove Theorem~\ref{thm:main}. All expectations are taken w.r.t. the randomness up to iteration $k$ unless otherwise noted.

\begin{proof}
\textbf{Step 1 (Smoothness application).}  
From Lemma~\ref{lem:smooth} with $g=\eta_k g_k$,
\begin{align}
f(w_{k+1})-f(w_k)
&\le -\eta_k\inner{\nabla f(w_k)}{g_k}+\frac{L}{2}\eta_k^{2}\|g_k\|^{2}. \tag{1}
\end{align}
[Step 1]

\textbf{Step 2 (Insert unbiasedness).}  
Take conditional expectation given $w_k$ and use $\E[g_k\mid w_k]=\nabla f(w_k)$:
\begin{align}
\E\!\big[f(w_{k+1})-f(w_k)\mid w_k\big]
&\le -\eta_k\|\nabla f(w_k)\|^{2}
+\frac{L}{2}\eta_k^{2}\E\!\big[\|g_k\|^{2}\mid w_k\big]. \tag{2}
\end{align}
[Step 2]

\textbf{Step 3 (Bound the second moment).}  
Decompose $\|g_k\|^{2}=\|\nabla f(w_k)\|^{2}+\|g_k-\nabla f(w_k)\|^{2}+2\inner{\nabla f(w_k)}{g_k-\nabla f(w_k)}$. The cross term vanishes after taking expectation, yielding
\[
\E\!\big[\|g_k\|^{2}\mid w_k\big]=\|\nabla f(w_k)\|^{2}+\E\!\big[\|g_k-\nabla f(w_k)\|^{2}\mid w_k\big]
\le \|\nabla f(w_k)\|^{2}+\sigma^{2}, \tag{3}
\]
by Assumption~\ref{ass:variance}.  
[Step 3]

\textbf{Step 4 (Insert (3) into (2)).}
\begin{align}
\E\!\big[f(w_{k+1})-f(w_k)\mid w_k\big]
&\le -\eta_k\|\nabla f(w_k)\|^{2}
+\frac{L}{2}\eta_k^{2}\big(\|\nabla f(w_k)\|^{2}+\sigma^{2}\big). \tag{4}
\end{align}
[Step 4]

\textbf{Step 5 (Choose $\eta_k$).}  
Recall $\eta_k=\dfrac{\alpha}{\sqrt{G_k+\epsilon}}$ and $G_k=\sum_{i=0}^{k} g_i^{2}$. Since $g_i^{2}\ge0$, we have $G_k\ge g_k^{2}$. Hence
\[
\eta_k^{2}\,g_k^{2}= \frac{\alpha^{2}g_k^{2}}{G_k+\epsilon}\le \alpha^{2}. \tag{5}
\]
[Step 5]

\textbf{Step 6 (Upper bound on the decrease).}  
Using $0\le \eta_k\le\alpha/\sqrt{\epsilon}$ and dropping the non‑negative term $\frac{L}{2}\eta_k^{2}\|\nabla f(w_k)\|^{2}$ from (4) gives a simpler inequality:
\begin{align}
\E\!\big[f(w_{k+1})-f(w_k)\mid w_k\big]
&\le -\frac{\alpha}{\sqrt{G_k+\epsilon}}\|\nabla f(w_k)\|^{2}
+\frac{L\alpha^{2}\sigma^{2}}{2\,(G_k+\epsilon)}. \tag{6}
\end{align}
[Step 6]

\textbf{Step 7 (Take full expectation).}  
Denote $\Delta_k:=\E[f(w_k)]-f^{\star}$. Taking total expectation of (6) and noting $\E[G_k]=\sum_{i=0}^{k}\E[g_i^{2}]$, we obtain
\begin{align}
\Delta_{k+1}-\Delta_{k}
&\le -\E\!\Big[\frac{\alpha}{\sqrt{G_k+\epsilon}}\|\nabla f(w_k)\|^{2}\Big]
+\frac{L\alpha^{2}\sigma^{2}}{2}\,\E\!\Big[\frac{1}{G_k+\epsilon}\Big]. \tag{7}
\end{align}
[Step 7]

\textbf{Step 8 (Summation over $k$).}  
Sum (7) from $k=0$ to $T-1$ and telescope the left–hand side:
\begin{align}
\Delta_{T}-\Delta_{0}
&\le -\alpha\sum_{k=0}^{T-1}\E\!\Big[\frac{\|\nabla f(w_k)\|^{2}}{\sqrt{G_k+\epsilon}}\Big]
+\frac{L\alpha^{2}\sigma^{2}}{2}\sum_{k=0}^{T-1}\E\!\Big[\frac{1}{G_k+\epsilon}\Big]. \tag{8}
\end{align}
Since $\Delta_{T}\ge0$, we can drop it and rearrange:
\[
\alpha\sum_{k=0}^{T-1}\E\!\Big[\frac{\|\nabla f(w_k)\|^{2}}{\sqrt{G_k+\epsilon}}\Big]
\le \Delta_{0}
+\frac{L\alpha^{2}\sigma^{2}}{2}\sum_{k=0}^{T-1}\E\!\Big[\frac{1}{G_k+\epsilon}\Big]. \tag{9}
\]
[Step 8]

\textbf{Step 9 (Bounding the variance term).}  
Using Lemma~\ref{lem:stepbound} with $a_k:=\E[g_k^{2}]\le G_{\max}$,
\[
\sum_{k=0}^{T-1}\E\!\Big[\frac{1}{G_k+\epsilon}\Big]
\le \frac{1}{\epsilon}\sum_{k=0}^{T-1}\E\!\Big[\frac{g_k^{2}}{G_k+\epsilon}\Big]
\le \frac{2}{\sqrt{\epsilon}}\sqrt{\sum_{k=0}^{T-1}\E[g_k^{2}]+\epsilon}
\le \frac{2\sqrt{T\,G_{\max}}}{\sqrt{\epsilon}}. \tag{10}
\]
A cruder but cleaner bound suffices for the theorem: since $G_k\ge g_k^{2}$,
\[
\frac{1}{G_k+\epsilon}\le\frac{1}{g_k^{2}+\epsilon}\le\frac{1}{\epsilon},
\]
hence
\[
\sum_{k=0}^{T-1}\E\!\Big[\frac{1}{G_k+\epsilon}\Big]\le\frac{T}{\epsilon}. \tag{10'}
\]
We will use (10') to keep constants simple.  
[Step 9]

\textbf{Step 10 (Relating the weighted sum to the plain sum).}  
Because $\sqrt{G_k+\epsilon}\le\sqrt{G_{T-1}+\epsilon}$ for all $k$, we have
\[
\frac{1}{\sqrt{G_{T-1}+\epsilon}}\sum_{k=0}^{T-1}\E\!\big[\|\nabla f(w_k)\|^{2}\big]
\le \sum_{k=0}^{T-1}\E\!\Big[\frac{\|\nabla f(w_k)\|^{2}}{\sqrt{G_k+\epsilon}}\Big]. \tag{11}
\]
[Step 10]

\textbf{Step 11 (Combine (9), (10') and (11)).}  
Insert (10') into (9) and then use (11):
\begin{align}
\frac{\alpha}{\sqrt{G_{T-1}+\epsilon}}\sum_{k=0}^{T-1}\E\!\big[\|\nabla f(w_k)\|^{2}\big]
&\le \Delta_{0}+ \frac{L\alpha^{2}\sigma^{2}}{2}\cdot\frac{T}{\epsilon}. \tag{12}
\end{align}
[Step 11]

\textbf{Step 12 (Lower bound on $\sqrt{G_{T-1}+\epsilon}$).}  
Since each $g_k^{2}\le G_{\max}$, we have $G_{T-1}\le T G_{\max}$. Consequently,
\[
\sqrt{G_{T-1}+\epsilon}\le \sqrt{T G_{\max}+\epsilon}\le \sqrt{T G_{\max}}+\sqrt{\epsilon}
\le 2\sqrt{T G_{\max}} \quad (\text{for }T\ge1,\;\epsilon\le TG_{\max}). \tag{13}
\]
Thus $1/\sqrt{G_{T-1}+\epsilon}\ge \frac{1}{2\sqrt{T G_{\max}}}$.  
[Step 12]

\textbf{Step 13 (Derive the average gradient norm bound).}  
Plug (13) into (12) and divide by $T$:
\begin{align}
\frac{1}{T}\sum_{k=0}^{T-1}\E\!\big[\|\nabla f(w_k)\|^{2}\big]
&\le \frac{2\sqrt{T G_{\max}}}{\alpha T}\Delta_{0}
+\frac{L\alpha\sigma^{2}}{\epsilon}\frac{1}{\sqrt{T G_{\max}}}. \tag{14}
\end{align}
Choosing $\epsilon=G_{\max}$ (a common practical choice) simplifies (14) to
\[
\frac{1}{T}\sum_{k=0}^{T-1}\E\!\big[\|\nabla f(w_k)\|^{2}\big]
\le \frac{2\Delta_{0}}{\alpha\sqrt{T}}+\frac{L\alpha\sigma^{2}}{\sqrt{T}}.
\]
Multiplying the first term by $1$ and the second by $1$ yields the bound stated in Theorem~\ref{thm:main}. ∎
\end{proof}

\section*{Rate Derivation (With Constants)}
From Theorem~\ref{thm:main} we extract the explicit constant:
\[
C_{1}= \frac{2\big(f(w_0)-f^{\star}\big)}{\alpha},\qquad
C_{2}=2L\alpha\sigma^{2}.
\]
Thus
\[
\frac{1}{T}\sum_{k=0}^{T-1}\E\!\big[\|\nabla f(w_k)\|^{2}\big]\le \frac{C_{1}+C_{2}}{\sqrt{T}}.
\]
The bound decays as $O(T^{-1/2})$ irrespective of convexity.

\section*{Special Cases}
\begin{enumerate}
\item \textbf{Deterministic gradients ($\sigma=0$).}  
The variance term disappears and we obtain
\[
\frac{1}{T}\sum_{k=0}^{T-1}\|\nabla f(w_k)\|^{2}\le \frac{2(f(w_0)-f^{\star})}{\alpha\sqrt{T}}.
\]
\item \textbf{Convex $f$.}  
When $f$ is convex, the same bound holds, and by convexity one can further translate the average gradient norm bound into a function‑value bound $O(1/\sqrt{T})$.
\item \textbf{Large base step size $\alpha$.}  
Increasing $\alpha$ improves the first term (faster descent) but linearly inflates the variance term $C_{2}$. The optimal $\alpha$ balancing both is $\alpha^{\star}=\sqrt{\frac{2(f(w_0)-f^{\star})}{L\sigma^{2}}}$, yielding a bound $\displaystyle O\big(\sigma\sqrt{L(f(w_0)-f^{\star})/T}\big)$.
\end{enumerate}

\end{document}